% Corrected French for printed p.151.
% No source correction admitted on this page.
\SourcePage{151}

\LayerTarget{tohoku:construction:resolving-functor-spectral-sequences:p2}
\noindent spectrales d'hyperhomologie (2.4.2) de $F$ par rapport au complexe
$K$.

\LayerTarget{tohoku:prop:2.5.3}
\noindent\textsc{Proposition 2.5.3.} --- \textit{Les isomorphismes (2.5.3)
proviennent d'isomorphismes fonctoriels de la première (resp. la deuxième)
suite spectrale du bicomplexe $\mathbf F(K)$ avec la première (resp. la
deuxième) suite spectrale d'hyperhomologie du foncteur $F$ par rapport au
complexe $K$.}

\LayerTarget{tohoku:prop:2.5.3:proof:p1}
Ces isomorphismes vont être explicités dans la démonstration ci-dessous. Soit
$C(K)=(C^{p,q}(K))_{p,q}$ une \og résolution injective\fg{} de $K$ au sens de
[6, Chap. XVII] ; considérant $K$ comme un bicomplexe à deuxième degré nul, on
a un homomorphisme d'augmentation $K\longrightarrow C(K)$. Considérons le
bicomplexe $M(K)=\mathbf F(C(K))$ donné par
\[
M^{p,q}(K)=\sum_{q'+q''=q}\mathbf F^{q''}(C^{p,q'}(K)).
\]
$M(K)$ n'est pas déterminé de façon unique par $K$, cependant les deux suites
spectrales de $M(K)$ et leur aboutissement commun le sont (puisque $C(K)$ est
déterminé à une équivalence d'homotopie de bicomplexes près), ce sont des
foncteurs spectraux bien déterminés en la variable $K$. Posons
$L(K)=F(C(K))$ ; les homomorphismes d'augmentation
$K\longrightarrow C(K)$ et $F\longrightarrow\mathbf F$ définissent des
homomorphismes de bicomplexes
\[
L(K)=F(C(K))\longrightarrow M(K)=\mathbf F(C(K))
\longleftarrow\mathbf F(K),
\]
d'où des homomorphismes \textit{fonctoriels} pour les suites spectrales
correspondantes (homomorphismes indépendants du choix de $C(K)$) :
\LayerTarget{tohoku:eq:2.5.4}
\[
\tag{2.5.4}
\begin{aligned}
\mathrm I L(K)&\longrightarrow\mathrm I M(K)\longleftarrow\mathrm I\mathbf F(K),\\
\mathrm{II}L(K)&\longrightarrow\mathrm{II}M(K)\longleftarrow\mathrm{II}\mathbf F(K),
\end{aligned}
\]
définissant les mêmes homomorphismes pour les aboutissements
\[
HL(K)\longrightarrow HM(K)\longleftarrow H\mathbf F(K).
\]
Pour les termes initiaux, les homomorphismes (2.5.4) donnent :
\LayerTarget{tohoku:eq:2.5.5}
\[
\tag{2.5.5}
\begin{aligned}
{}_{\mathrm I}E_2^{p,q}L(K)&\longrightarrow{}_{\mathrm I}E_2^{p,q}M(K)
\longleftarrow{}_{\mathrm I}E_2^{p,q}\mathbf F(K),\\
{}_{\mathrm{II}}E_2^{p,q}L(K)&\longrightarrow{}_{\mathrm{II}}E_2^{p,q}M(K)
\longleftarrow{}_{\mathrm{II}}E_2^{p,q}\mathbf F(K).
\end{aligned}
\]
Nous allons prouver que les homomorphismes (2.5.5) sont des isomorphismes, et
que les isomorphismes correspondants entre les termes extrêmes de (2.5.5)
(qui sont les termes initiaux des suites spectrales d'hyperhomologie de $F$
pour $K$, resp. du bicomplexe $\mathbf F(K)$) sont ceux résultant des formules
(2.5.3) ; il en résultera que les homomorphismes (2.5.4) sont des
isomorphismes, d'où des isomorphismes entre les termes extrêmes, ce seront
évidemment les isomorphismes cherchés. Tout revient donc à prouver que les
termes médians des lignes (2.5.5) ont respectivement la forme
$H^p(R^q(F(K)))$ et $R^pF(H^q(K))$ (la vérification que les homomorphismes
dans (2.5.5) sont bien les isomorphismes naturels étant alors purement
mécanique, et d'ailleurs implicitement contenue dans la démonstration qui
suit).

On va démontrer $H_{\mathrm{II}}(M)^{p,q}=R^qF(K^p)$ (d'où résultera aussitôt
$ {}_{\mathrm I}E_2^{p,q}(M)=H^p(R^qF(K))$). Or pour $p$ fixé,
$M^{p*}=(M^{p,q})_q$ est le complexe simple associé au bicomplexe
\[
(\mathbf F^{q''}(C^{p,q'}(K)))_{q',q''}=\mathbf F(C(K^p)),
\]
où $C(K^p)$ désigne le complexe $C^{p*}(K)$,

\TohokuContinuation{152}

